\subsubsection{String Manipulation}

Positions are 1-based, and every function returns a new value rather than
mutating in place. The recurring uses are normalizing a dirty join key, parsing
a delimited field, and folding a group of strings into one.

\paragraph{Concatenation.} The operator \texttt{a || b} propagates
\texttt{NULL}: if either side is \texttt{NULL} the whole expression is
\texttt{NULL}. \texttt{CONCAT(a, b, ...)} treats \texttt{NULL} as an empty
string. \texttt{CONCAT\_WS(sep, ...)} takes the separator first, joins the rest
with it, and skips \texttt{NULL} arguments without leaving a doubled separator.
Choosing \texttt{||} for user-facing text is a common way to silently blank out
a whole column.

\paragraph{Case and whitespace.} \texttt{UPPER} and \texttt{LOWER} are the
standard folds, \texttt{INITCAP} capitalizes the first letter of each word and
lowercases the rest. \texttt{TRIM} strips characters from the ends only, spaces
by default, with \texttt{LEADING}, \texttt{TRAILING}, or \texttt{BOTH} to pick
a side. Interior whitespace survives, so collapsing double spaces needs
\texttt{REGEXP\_REPLACE(s, '\textbackslash s+', ' ', 'g')}.
\texttt{LPAD(s, n, fill)} and \texttt{RPAD} pad to a fixed width and truncate
when the input is already longer than \texttt{n}.

\vspace{2ex}

\begin{lstlisting}[style=sql, caption={Normalizing a join key}]
SELECT
    LOWER(TRIM(email))                      AS email_key,
    SPLIT_PART(LOWER(TRIM(email)), '@', 2)  AS domain,
    LPAD(account_id::text, 8, '0')          AS padded_id,
    CONCAT_WS(', ', last_name, first_name)  AS sort_name
FROM users;
\end{lstlisting}

\newpage

\paragraph{Extraction and search.} \texttt{SUBSTRING(s FROM a FOR b)} (or
\texttt{SUBSTR(s, a, b)}) takes \texttt{b} characters starting at position
\texttt{a}. \texttt{LEFT(s, n)} and \texttt{RIGHT(s, n)} take from the ends,
and a negative \texttt{n} drops that many characters from the opposite end.
\texttt{STRPOS(s, sub)} returns the 1-based index of the first occurrence and
\texttt{0} when the substring is absent, so a found/not-found test compares
against \texttt{0}. \texttt{SPLIT\_PART(s, delim, n)} pulls the \texttt{n}th
field of a delimited string and returns an empty string when that field does
not exist, worth wrapping in \texttt{NULLIF(..., '')}.

\begin{lstlisting}[style=sql, caption={Parsing a delimited name field}]
SELECT full_name,
       SPLIT_PART(full_name, ' ', 1)              AS first_name,
       NULLIF(SPLIT_PART(full_name, ' ', 2), '')  AS last_name,
       LENGTH(full_name)
         - LENGTH(REPLACE(full_name, ' ', ''))    AS space_count
FROM people;
\end{lstlisting}

\vspace{-1ex}

The last expression is the occurrence-counting idiom: remove every copy of the
target and measure how much the string shrank. For a needle longer than one
character, divide the difference by \texttt{LENGTH(needle)}.

\paragraph{Replacement.} \texttt{REPLACE(s, from, to)} substitutes every
literal occurrence. \texttt{TRANSLATE(s, from, to)} maps character to character
by position and deletes any character in \texttt{from} that has no counterpart
in \texttt{to}, a fast way to strip a known character set.
\texttt{REGEXP\_REPLACE(s, pattern, repl, flags)} handles the general case and
replaces only the \emph{first} match unless the \texttt{`g'} flag is passed, a
frequent source of half-cleaned output.

\paragraph{Pattern matching.} \texttt{LIKE} uses \texttt{\%} for any run of
characters and \texttt{\_} for exactly one, and it is case sensitive.
\texttt{ILIKE} is the case-insensitive form. For real expressiveness the POSIX
operators apply: \texttt{\textasciitilde} matches, \texttt{\textasciitilde*}
matches case-insensitively, and \texttt{!\textasciitilde} negates. Anchors
\texttt{\^{}} and \texttt{\$} matter, since an unanchored pattern matches
anywhere in the string.

\begin{lstlisting}[style=sql, caption={Regex cleaning and validation}]
SELECT phone,
       REGEXP_REPLACE(phone, '[^0-9]', '', 'g')  AS digits,
       phone ~ '^\+?[0-9]{10,15}$'               AS looks_valid,
       email ILIKE '%@gmail.com'                 AS is_gmail
FROM contacts;
\end{lstlisting}

\vspace{-1ex}

\paragraph{Splitting and aggregating.} \texttt{STRING\_AGG(expr, delim ORDER BY
...)} collapses a group into one delimited string and is the string analogue of
\texttt{SUM}. The \texttt{ORDER BY} lives inside the call, so the output order
is deterministic. The inverse direction expands one row into many:
\texttt{UNNEST(STRING\_TO\_ARRAY(s, delim))} or
\texttt{REGEXP\_SPLIT\_TO\_TABLE(s, pattern)} in the \texttt{FROM} clause.

\begin{lstlisting}[style=sql, caption={Collapse a group, then expand it back}]
SELECT department,
       STRING_AGG(name, ', ' ORDER BY name) AS members
FROM employees
GROUP BY department;
-- one row per tag stored in a comma-separated column
SELECT p.id, TRIM(tag) AS tag
FROM products p,
     UNNEST(STRING_TO_ARRAY(p.tags, ',')) AS tag;
\end{lstlisting}
